\documentclass[10pt,a4paper]{amsart}
\usepackage{amsmath,amssymb,amsthm,graphicx}
\usepackage{verbatim}
\newenvironment{reference}{\comment}{\endcomment}
%\newenvironment{reference}{}{}
\newenvironment{slogan}{\comment}{\endcomment}
\newenvironment{history}{\comment}{\endcomment}
% For commutative diagrams we use Xy-pic
\usepackage[all]{xy}

% We use 2cell for 2-commutative diagrams.
\xyoption{2cell}
% Theorem environments.
%
\theoremstyle{plain}
\newtheorem{theorem}[subsection]{Theorem}
\newtheorem{proposition}[subsection]{Proposition}
\newtheorem{lemma}[subsection]{Lemma}

\theoremstyle{definition}
\newtheorem{definition}[subsection]{Definition}
\newtheorem{example}[subsection]{Example}
\newtheorem{exercise}[subsection]{Exercise}
\newtheorem{situation}[subsection]{Situation}

\theoremstyle{remark}
\newtheorem{remark}[subsection]{Remark}
\newtheorem{remarks}[subsection]{Remarks}

\numberwithin{equation}{subsection}

% Macros
%
\def\lim{\mathop{\mathrm{lim}}\nolimits}
\def\colim{\mathop{\mathrm{colim}}\nolimits}
\def\Spec{\mathop{\mathrm{Spec}}}
\def\Hom{\mathop{\mathrm{Hom}}\nolimits}
\def\Ext{\mathop{\mathrm{Ext}}\nolimits}
\def\SheafHom{\mathop{\mathcal{H}\!\mathit{om}}\nolimits}
\def\SheafExt{\mathop{\mathcal{E}\!\mathit{xt}}\nolimits}
\def\Sch{\mathit{Sch}}
\def\Mor{\mathop{\mathrm{Mor}}\nolimits}
\def\Ob{\mathop{\mathrm{Ob}}\nolimits}
\def\Sh{\mathop{\mathit{Sh}}\nolimits}
\def\NL{\mathop{N\!L}\nolimits}
\def\CH{\mathop{\mathrm{CH}}\nolimits}
\def\proetale{{pro\text{-}\acute{e}tale}}
\def\etale{{\acute{e}tale}}
\def\QCoh{\mathit{QCoh}}
\def\Ker{\mathop{\mathrm{Ker}}}
\def\Im{\mathop{\mathrm{Im}}}
\def\Coker{\mathop{\mathrm{Coker}}}
\def\Coim{\mathop{\mathrm{Coim}}}

% Boxtimes
%
\DeclareMathSymbol{\boxtimes}{\mathbin}{AMSa}{"02}

%
% Macros for moduli stacks/spaces
%
\def\QCohstack{\mathcal{QC}\!\mathit{oh}}
\def\Cohstack{\mathcal{C}\!\mathit{oh}}
\def\Spacesstack{\mathcal{S}\!\mathit{paces}}
\def\Quotfunctor{\mathrm{Quot}}
\def\Hilbfunctor{\mathrm{Hilb}}
\def\Curvesstack{\mathcal{C}\!\mathit{urves}}
\def\Polarizedstack{\mathcal{P}\!\mathit{olarized}}
\def\Complexesstack{\mathcal{C}\!\mathit{omplexes}}
% \Pic is the operator that assigns to X its picard group, usage \Pic(X)
% \Picardstack_{X/B} denotes the Picard stack of X over B
% \Picardfunctor_{X/B} denotes the Picard functor of X over B
\def\Pic{\mathop{\mathrm{Pic}}\nolimits}
\def\Picardstack{\mathcal{P}\!\mathit{ic}}
\def\Picardfunctor{\mathrm{Pic}}
\def\Deformationcategory{\mathcal{D}\!\mathit{ef}}
\usepackage{cleveref}
\crefname{originalequation}{original equation}{original equations}
\Crefname{originalequation}{Original equation}{Original equations}
\makeatletter
\newlabel{equation-base-change}{{equation-base-change}{source}{}{sourcepointer-1}{}}
\newlabel{equation-base-change@cref}{{[originalequation][0][]equation-base-change}{[1][0][]source}{}{}{}}
\makeatother
\crefname{originalreference}{original reference}{original references}
\Crefname{originalreference}{Original reference}{Original references}
\makeatletter
\newlabel{more-algebra-equation-base-change-RHom}{{more-algebra-equation-base-change-RHom}{source}{}{sourcepointer-2}{}}
\newlabel{more-algebra-equation-base-change-RHom@cref}{{[originalreference][0][]more-algebra-equation-base-change-RHom}{[1][0][]source}{}{}{}}
\makeatother
\crefname{originalreference}{original reference}{original references}
\Crefname{originalreference}{Original reference}{Original references}
\makeatletter
\newlabel{more-algebra-lemma-base-change-RHom}{{more-algebra-lemma-base-change-RHom}{source}{}{sourcepointer-3}{}}
\newlabel{more-algebra-lemma-base-change-RHom@cref}{{[originalreference][0][]more-algebra-lemma-base-change-RHom}{[1][0][]source}{}{}{}}
\makeatother
\crefname{originalreference}{original reference}{original references}
\Crefname{originalreference}{Original reference}{Original references}
\makeatletter
\newlabel{more-algebra-lemma-internal-hom-diagonal-better}{{more-algebra-lemma-internal-hom-diagonal-better}{source}{}{sourcepointer-4}{}}
\newlabel{more-algebra-lemma-internal-hom-diagonal-better@cref}{{[originalreference][0][]more-algebra-lemma-internal-hom-diagonal-better}{[1][0][]source}{}{}{}}
\makeatother
\crefname{originalreference}{original reference}{original references}
\Crefname{originalreference}{Original reference}{Original references}
\makeatletter
\newlabel{more-algebra-lemma-upgrade-adjoint-tensor-RHom}{{more-algebra-lemma-upgrade-adjoint-tensor-RHom}{source}{}{sourcepointer-5}{}}
\newlabel{more-algebra-lemma-upgrade-adjoint-tensor-RHom@cref}{{[originalreference][0][]more-algebra-lemma-upgrade-adjoint-tensor-RHom}{[1][0][]source}{}{}{}}
\makeatother
\title{Testing the base change map after pushing forward along the base change}
\author{The Stacks Project Authors}
\date{stacks-project:dualizing.tex:local-snapshot-20261012; GNU FDL 1.2 or later}
\begin{document}
\maketitle
\noindent\textit{Editorial note.} All delivered mathematical text is the author's own, copied from the pinned Stacks Project chapter source sources/stacks-dualizing/pinned-original-source.tex; no statement or proof step was rewritten, reordered or removed.\par\medskip
\noindent\textit{Editorial note.} Preamble provenance: the Stacks Project keeps one single shared preamble for every chapter; this card takes it directly from the staged sources/stacks-dualizing/preamble.tex as three verbatim spans (lines 8-12, 14-18 and 154-217), so no derived file was written and nothing was added or redefined. The wrapper class and the base packages are supplied by the assembler.\par\medskip
\noindent\textit{Editorial note.} Bibliography provenance: neither the counted statement nor its proof contains a citation (this lane rescans the whole span before writing the card), so the card carries no bibliography and no entry was reconstructed.\par\medskip
\noindent\textit{Editorial note.} Selection provenance: this card was opened under the file-based control contract (CONTROL.md), with the mechanical pre-card gate run on its statement and proof interval before the card was written.\par\medskip
\noindent\textit{Allowed standard prerequisites (editorial).}\begin{itemize}
\item The possibility of testing an isomorphism after applying the pushforward along the base changed ring map. Invoked at the start of the counted proof (source lines 2709-2710). Reduces the statement to a comparison after pushing forward.
\item The base change map for the derived internal hom and its identification with the map obtained after pushing forward. Invoked inside the counted proof (source lines 2711-2717). Identifies the pushed forward map with the known base change map for the internal hom.
\item The unwinding of the two sides of that map along the adjunction between tensor and internal hom. Invoked at the end of the counted proof (source lines 2718-2722). Produces the explicit comparison whose isomorphism is the assertion.
\end{itemize}\medskip
\ifdefined\hypertarget\hypertarget{sourcepointer-1}{}\fi
\noindent\textit{Source reference.} Original equation equation-base-change: sources/stacks-dualizing/pinned-original-source.tex, source lines 2669--2669. Navigation pointer only: the counted text refers to the Stacks item equation-base-change, which is an existing labeled result of this same chapter and is not part of this card. The pointer carries the target's own label as its number and its definition line as its span; its content is not reproduced here.\par\medskip
\ifdefined\hypertarget\hypertarget{sourcepointer-2}{}\fi
\noindent\textit{Source reference.} Original reference more-algebra-equation-base-change-RHom: sources/stacks-dualizing/pinned-original-source.tex, source lines 2717--2717. Navigation pointer only: the counted text refers to the Stacks item more-algebra-equation-base-change-RHom, which belongs to another chapter of the project and is not part of this card. The pointer carries the target's own label as its number because this chapter's staged snapshot does not contain the other chapter's tag table; its content is not reproduced here.\par\medskip
\ifdefined\hypertarget\hypertarget{sourcepointer-3}{}\fi
\noindent\textit{Source reference.} Original reference more-algebra-lemma-base-change-RHom: sources/stacks-dualizing/pinned-original-source.tex, source lines 2719--2719. Navigation pointer only: the counted text refers to the Stacks item more-algebra-lemma-base-change-RHom, which belongs to another chapter of the project and is not part of this card. The pointer carries the target's own label as its number because this chapter's staged snapshot does not contain the other chapter's tag table; its content is not reproduced here.\par\medskip
\ifdefined\hypertarget\hypertarget{sourcepointer-4}{}\fi
\noindent\textit{Source reference.} Original reference more-algebra-lemma-internal-hom-diagonal-better: sources/stacks-dualizing/pinned-original-source.tex, source lines 2705--2705. Navigation pointer only: the counted text refers to the Stacks item more-algebra-lemma-internal-hom-diagonal-better, which belongs to another chapter of the project and is not part of this card. The pointer carries the target's own label as its number because this chapter's staged snapshot does not contain the other chapter's tag table; its content is not reproduced here.\par\medskip
\ifdefined\hypertarget\hypertarget{sourcepointer-5}{}\fi
\noindent\textit{Source reference.} Original reference more-algebra-lemma-upgrade-adjoint-tensor-RHom: sources/stacks-dualizing/pinned-original-source.tex, source lines 2721--2721. Navigation pointer only: the counted text refers to the Stacks item more-algebra-lemma-upgrade-adjoint-tensor-RHom, which belongs to another chapter of the project and is not part of this card. The pointer carries the target's own label as its number because this chapter's staged snapshot does not contain the other chapter's tag table; its content is not reproduced here.\par\medskip
\begin{lemma}
\label{lemma-check-base-change-is-iso}
In the situation above, the map (\ref{equation-base-change})
is an isomorphism if and only if the map
$$
R\Hom_R(A, K) \otimes_R^\mathbf{L} R'
\longrightarrow
R\Hom_R(A, K \otimes_R^\mathbf{L} R')
$$
of More on Algebra, Lemma
\ref{more-algebra-lemma-internal-hom-diagonal-better} is an isomorphism.
\end{lemma}

\begin{proof}
To see that the map is an isomorphism, it suffices to prove it
is an isomorphism after applying $\varphi'_*$.
Applying the functor $\varphi'_*$ to (\ref{equation-base-change})
and using that $A' = A \otimes_R^\mathbf{L} R'$
we obtain the base change map
$R\Hom_R(A, K) \otimes_R^\mathbf{L} R' \to
R\Hom_{R'}(A \otimes_R^\mathbf{L} R', K \otimes_R^\mathbf{L} R')$
for derived hom of
More on Algebra, Equation (\ref{more-algebra-equation-base-change-RHom}).
Unwinding the left and right hand side exactly as in the proof of
More on Algebra, Lemma \ref{more-algebra-lemma-base-change-RHom}
and in particular using
More on Algebra, Lemma \ref{more-algebra-lemma-upgrade-adjoint-tensor-RHom}
gives the desired result.
\end{proof}

\end{document}
